\Cref{fig:1-challenges} connects three challenges to the mechanisms and
evidence used to address them.

First, the primary codebook has useful mathematical structure that a direct
GPU port can overlook. Each eight-bit subspace code indexes 256 distances.
Its Cartesian coordinates permit an exact 16+16 factorisation, but more
aggressive splitting adds lookups and arithmetic. The challenge is to exploit
separability without assuming that the smallest table gives the fastest
scan. Candidate-parallel word access and residual lane cooperation also need
different data mappings.

Second, local quotas can lose global winners. For disjoint groups $U_1,U_2$ and $c=4$, consider
\[
 \begin{aligned}
 \operatorname{Top}_4(U)&\subseteq U_1,\\
 \left|\operatorname{Top}_4(U)\cap
 \bigl(\operatorname{Top}_2(U_1)\cup\operatorname{Top}_2(U_2)\bigr)\right|
 &=2<4.
 \end{aligned}
\]
Groups need coverage before merging. Bounds reduce maintenance; scheduling controls parallelism. IVF routing is separate.

Third, reference agreement cannot fix reference error. A $k=1$ counterexample is
\[
 \mathcal R_\alpha=\mathcal R_{\rm ref}=\{a\},\quad
 \mathcal T_1=\{b\},\quad a\ne b
 \ \Longrightarrow\ 
 D_S(\alpha)=0,\quad \Recall@1=0
 \qquad(S=1).
\]
We test quality alongside calibration/reuse cost.

\par\medskip\noindent\begin{minipage}{\linewidth}
\centering\captionsetup{hypcap=false}
\resizebox{\linewidth}{!}{\begin{tikzpicture}[
 font=\footnotesize\linespread{1}\selectfont,
 box/.style={draw=ThesisNavy!65,rounded corners=2pt,align=center,
 text width=3.5cm,minimum height=1.20cm,inner sep=3pt},
 arr/.style={-{Latex[length=2mm]},draw=ThesisNavy!80,thick}]
\node[font=\bfseries] at (0,.80) {Challenge};
\node[font=\bfseries] at (4.6,.80) {Mechanism};
\node[font=\bfseries] at (9.2,.80) {Evidence};
\node[box] (a) at (0,0) {Primary codes\\Broad scan cost};
\node[box,fill=ThesisNavy!7] (b) at (4.6,0) {Cartesian tables\\Packed word scan};
\node[box] (c) at (9.2,0) {Table / load tests\\Score checks};
\node[box] (d) at (0,-1.45) {Uneven query work\\Global top-$c$};
\node[box,fill=ThesisNavy!7] (e) at (4.6,-1.45) {Grouped scan\\Shared bound\\Global merge};
\node[box] (f) at (9.2,-1.45) {Output checks\\Batch / skew tests};
\node[box] (g) at (0,-2.90) {Refinement effort\\Quality / reuse cost};
\node[box,fill=ThesisNavy!7] (h) at (4.6,-2.90) {Sampled agreement\\Workload budget};
\node[box] (i) at (9.2,-2.90) {Held-out quality\\Calibration cost};
\draw[arr] (a)--(b);\draw[arr] (b)--(c);
\draw[arr] (d)--(e);\draw[arr] (e)--(f);
\draw[arr] (g)--(h);\draw[arr] (h)--(i);
\end{tikzpicture}
}
\captionof{figure}[Problem, mechanism and evidence]{Problem, mechanism and evidence. The arrows organise
research questions; they do not imply that independently measured speedups
multiply or that an unlabelled agreement rule guarantees GT recall.}
\label{fig:1-challenges}
\end{minipage}\par\medskip
